\ncertchapter{5}{Kinetic Theory}\chlabel{ch:kinetic}
% ================================================================

Part~A of this volume treated matter in bulk: solids that stretch and fluids that push, described by
moduli, pressure and viscosity, with no mention of atoms. Part~B introduced temperature and heat, and
\chref{ch:thermo} turned them into laws. Throughout, a gas was a black box obeying $PV=nRT$, with an
internal energy that depended on temperature alone; heat went in, work came out, and nobody asked what
the gas is \emph{made of}. This chapter opens the box.

The idea is simple to state and was bitterly resisted for most of the nineteenth century: a gas is an
enormous crowd of tiny molecules flying about at random, colliding with each other and with the walls,
and obeying the ordinary mechanics of \chsref{ch:newton}{ch:collisions}. From that picture alone we shall
derive the pressure of a gas, discover what temperature \emph{is}, explain the gas laws that
\chref{ch:thermo} took as given, predict the heat capacities of gases from the shapes of their molecules,
and estimate how far a molecule travels before it hits another. Each derivation uses only tools already
built: momentum and impulse from Volume~I, and averages over a distribution from \chref{ch:toolkit}.

The limits deserve stating at the outset. Kinetic theory as developed here is a theory of the
\emph{ideal} gas: point-like molecules that interact only when they collide. It fails, and says that it
will fail, when a gas is dense or cold enough for the size of its molecules and the forces between them
to matter. And at one decisive point, the heat capacities at low temperature, classical mechanics itself
gives the wrong answer, which was one of the first signs that quantum physics was needed.

\begin{motivbox}{Why kinetic theory matters}
Kinetic theory connects the two halves of this volume: it \emph{derives} the gas laws and the meaning of
temperature from nothing but molecules moving under Newton's laws.
\begin{itemize}
\item \textbf{It explains what temperature actually is.} The temperature of a gas turns out to measure
the average kinetic energy of its molecules, $\tfrac32 k_BT$ per molecule. This is why gases exert
pressure (molecules drumming on the walls), why a hot gas pushes harder, and why absolute zero is a floor.
\item \textbf{It predicts the gas laws and specific heats from first principles.} Boyle's law, Charles's
law and $PV=nRT$ all follow from counting molecular impacts; the equipartition of energy then predicts
the specific heats of gases, and its failure at low temperature gave the first hints of quantum mechanics.
\item \textbf{It underlies real technology.} The mean free path decides how a vacuum pump, a thermos
flask or a semiconductor deposition chamber behaves; molecular speeds explain why light gases (hydrogen,
helium) escape a planet's atmosphere while heavier ones stay, and how uranium isotopes are separated by
effusion.
\end{itemize}
\end{motivbox}

% ================================================================
\section{The Molecular Picture and the Gas Laws}\label{sec:ktpicture}

Before any derivation, the model must be stated precisely, because every later result inherits its
assumptions and, with them, its limits.

\subsection{The postulates, stated honestly}

Kinetic theory rebuilds the gas laws from one radical premise: that a gas is a vast number of tiny
molecules in ceaseless random motion. Two counts are used side by side. The \textbf{number of molecules} is $N$, a pure
number; the \textbf{amount of substance} is $n$ moles, as in \chref{ch:thermal} (Section~\ref{sec:gaslaw}), where the mole was introduced. Recall that they are linked by Avogadro's number $N_A$, the number of molecules in one mole: $N=nN_A$. Each molecule has mass $m$ (in
kilograms), and a mole of them has the \textbf{molar mass} $M=N_Am$ (in \si{kg/mol}).

\begin{defbox}{Assumptions of kinetic theory (the ideal gas)}
\begin{enumerate}[label=\textbf{(\roman*)}]
\item A gas consists of a very large number $N$ of identical molecules in random motion.
\item The volume of the molecules themselves is negligible compared with the volume of the container
      (about $10^{-3}$ of it for air at room conditions; see below).
\item There are no forces between molecules except during collisions; between collisions each molecule
      moves in a straight line at constant speed (gravity is neglected).
\item Collisions, with each other and with the walls, are perfectly \emph{elastic} and of negligible
      duration.
\item The number $N$ is so large (about $10^{23}$) that averages are meaningful and fluctuations about
      them are far too small to observe.
\end{enumerate}
\end{defbox}

\why{Each assumption earns its keep, and each tells you where the theory will fail. Assumptions (ii) and
(iii) break down at high pressure or low temperature, when molecules are crowded together and move slowly
enough for their mutual attraction to matter. That is precisely when real gases stop obeying $PV=nRT$ and
begin to liquefy. A theory that predicts the conditions of its own breakdown is a good theory.}

How empty is a gas? Air at \SI{0}{\celsius} and atmospheric pressure holds about
$n_V=\SI{2.7e25}{}$ molecules per cubic metre (the \textbf{number density} $n_V=N/V$, in \si{m^{-3}};
we derive this value shortly). Each molecule therefore has a little cube of volume $1/n_V$ to itself, of
side $n_V^{-1/3}\approx\SI{3.3}{nm}$. A nitrogen or oxygen molecule is about $d\approx\SI{0.37}{nm}$
across, so the molecules are spaced about nine diameters apart, and they fill only
$n_V\cdot\tfrac{\pi}{6}d^3\approx7\times10^{-4}$ of the space. Assumption (ii) is good to a fraction of a
per cent for air at room conditions; it is the liquid, a thousand times denser, where it fails.

\subsection{The ideal-gas law, counted per molecule}

\chref{ch:thermal} (Section~\ref{sec:gaslaw}) wrote the ideal-gas law per mole, and \chref{ch:thermo} used it in that form throughout. Kinetic theory prefers to count molecules.

\begin{keybox}{The ideal-gas law, per mole and per molecule}
With $N=nN_A$,
\[
PV = nRT = Nk_BT, \qquad k_B = \frac{R}{N_A} = \SI{1.38e-23}{J/K}.
\]
Here $P$ is the pressure (\si{Pa}), $V$ the volume (\si{m^3}) and $T$ the absolute temperature (\si{K}).
$R=\SI{8.314}{J.mol^{-1}.K^{-1}}$ is the gas constant \emph{per mole}; $k_B$, Boltzmann's constant, is
the same constant \emph{per molecule}. Dividing by $V$ gives the most useful form for this chapter,
\[
P = n_Vk_BT, \qquad n_V=\frac{N}{V}.
\]
\end{keybox}

So far this is an experimental fact, summarising the gas laws of Boyle and Charles. The task of the next
two sections is to \emph{derive} it: after Section~\ref{sec:kttemperature} it becomes a theorem about
molecules, with temperature given a mechanical meaning. As a first use, the number density quoted above
follows at once: at $P=\SI{1.013e5}{Pa}$ and $T=\SI{273}{K}$,
$n_V=P/k_BT=\SI{2.69e25}{m^{-3}}$.

\why{Avogadro's number $N_A$ is the exchange rate between the two worlds this chapter connects. The
laboratory measures $P$, $V$, $T$ and $R$, all \emph{bulk} quantities. Kinetic theory speaks of single
molecules: their mass $m=M/N_A$, their energy, their count $N=nN_A$. Every time we cross from one
description to the other we multiply or divide by $N_A$, and the relation $k_B=R/N_A$ is that crossing
written once and for all: the per-molecule constant is the per-mole constant shared out among Avogadro's
number of molecules. Why the same $N_A$ works for every gas, and how it was measured, is taken up in
Section~\ref{sec:kttemperature}, once temperature has its kinetic meaning.}

% ================================================================
\section{Pressure from Molecular Impacts}\label{sec:ktpressure}

The central derivation of the chapter shows that pressure is nothing but the steady drumbeat of molecular
collisions on the walls. It needs two results from Volume~I. First, force is the rate of change of
momentum, $\vF=d\vp/dt$, so a force acting for a short time delivers an impulse equal to the change in
momentum it produces (\chref{ch:newton}). Second, when a ball bounces elastically off a smooth, massive
wall, the component of its velocity normal to the wall is reversed and the component along the wall is
unchanged (\chref{ch:collisions}): the wall is so heavy that it takes up momentum without taking up any
appreciable kinetic energy. Figure~\ref{fig:ktwall} sets up the calculation.

\begin{figure}[htbp]
\centering
\tikzfig{kinetic_wall}
\figcaption{fig:ktwall}{Left: one molecule in a cube of side $\ell$. Only its $x$-component of velocity carries it towards the shaded wall of area $\ell^2$, and it returns to that wall after a round trip of $2\ell$, so it hits at a steady rate. Right: one elastic hit seen from above. The component normal to the wall reverses, $v_x\to-v_x$, while $v_y$ is untouched, so each hit hands the wall momentum $2mv_x$.}
\end{figure}

\begin{derivbox}{The pressure of an ideal gas}
Take a cube of side $\ell$ containing $N$ molecules, each of mass $m$, and fix attention on the wall
perpendicular to the $x$-axis at $x=\ell$ (Figure~\ref{fig:ktwall}).

\stepn{1} \textbf{One molecule, one hit.} A molecule with velocity component $v_x>0$ hits the wall and
rebounds elastically, $v_x\to-v_x$. Its momentum changes by $\Delta p_x=-2mv_x$, so by the third law it
delivers $+2mv_x$ to the wall.

\stepn{2} \textbf{How often?} Between successive hits on the \emph{same} wall the molecule must cross to the
far wall and back, a distance $2\ell$ in the $x$-direction at speed $|v_x|$. That takes a time $2\ell/|v_x|$,
so it strikes the wall $|v_x|/2\ell$ times per second.

\stepn{3} \textbf{Average force from one molecule} $=$ momentum delivered per second:
\[
F_1 = 2m|v_x| \times \frac{|v_x|}{2\ell} = \frac{m v_x^2}{\ell}.
\]
The absolute values merge into a plain square, so the sign of $v_x$ no longer matters.

\stepn{4} \textbf{All $N$ molecules.} Add the contributions and write the sum as $N$ times the
\emph{average} of $v_x^2$ (a bar denotes the average over all the molecules):
\[
F = \frac{m}{\ell}\sum_{i=1}^{N} v_{x,i}^2 = \frac{mN\overline{v_x^2}}{\ell}.
\]
A real molecule also moves along $y$ and $z$, but only its $x$-motion carries momentum to \emph{this}
wall, and its $y$- and $z$-motion neither speed up nor slow down its return in $x$; so the count in
Step~2 stands.

\stepn{5} \textbf{Isotropy.} No direction is special, so
$\overline{v_x^2} = \overline{v_y^2} = \overline{v_z^2}$. Since $v^2 = v_x^2+v_y^2+v_z^2$ for every
molecule, the averages obey $\overline{v^2}=\overline{v_x^2}+\overline{v_y^2}+\overline{v_z^2}$, and so
\[
\overline{v_x^2} = \tfrac13\overline{v^2}.
\]

\stepn{6} \textbf{Pressure.} Divide the force by the wall's area $\ell^2$, and write $V=\ell^3$ and
$\rho = Nm/V$ for the density of the gas:
\[
P = \frac{F}{\ell^2} = \frac{mN\overline{v^2}}{3\ell^3}
\quad\Longrightarrow\quad
\boxed{\;P = \frac{1}{3}\,\frac{Nm}{V}\,\overline{v^2} = \frac{1}{3}\rho\,\overline{v^2}\;}
\]

\stepn{7} \textbf{What was set aside.} Collisions between molecules were ignored. They do not change the
result: momentum and kinetic energy are conserved in every collision, and in equilibrium the gas as a
whole keeps the same distribution of velocities, so the stream of molecules arriving at the wall has the
same mean value of $v_x^2$ as before. Nor does the shape of the container matter: the pressure is a local
property, and any wall can be cut into small flat patches, each of which sees the same molecular rain.
\end{derivbox}

\inwords{Pressure is a rain of impacts. It grows with the number of molecules per unit volume (more
drummers), with the molecular mass (heavier drumsticks), and with the mean-square speed: a faster molecule
hits both harder \emph{and} more often, which is why speed enters squared.}

Two rewritings of the result are worth having. Since the mean translational kinetic energy of a molecule
is $\overline K=\tfrac12m\overline{v^2}$ (we write $K$ for kinetic energy, as in Volume~I),
\[
P=\tfrac23\,n_V\,\overline K ,
\]
so pressure is two-thirds of the translational kinetic energy per unit volume: a pressure in pascals is
also an energy density in joules per cubic metre. And because each species of molecule in a mixture
drums on the wall independently of the others, the pressures they would exert alone simply add:
\[
P = P_1 + P_2 + \cdots, \qquad P_i = n_{V,i}\,k_BT .
\]
This is \textbf{Dalton's law of partial pressures}, which here costs nothing: it is the sum in Step~4
split by species.

\begin{exbox}{How fast are the molecules of air?}\label{ex:ktairspeed}
The pressure formula can be run backwards. Air at sea level and \SI{0}{\celsius} has
$P=\SI{1.013e5}{Pa}$ and density $\rho=\SI{1.29}{kg/m^3}$. From $P=\tfrac13\rho\overline{v^2}$,
\[
\sqrt{\overline{v^2}}=\sqrt{\frac{3P}{\rho}}=\sqrt{\frac{3(1.013\times10^5)}{1.29}}\approx\SI{485}{m/s}.
\]
Two measurable bulk quantities give a molecular speed, about \SI{1750}{km/h}, without knowing the mass or
the number of molecules. The square root of the mean square speed is called the rms speed; it is
studied properly in the next section. \ans{about \SI{485}{m/s}}
\end{exbox}

% ================================================================
\section{Temperature and Molecular Speeds}\label{sec:kttemperature}

The pressure derivation gave $PV = \tfrac13 Nm\overline{v^2}$, a purely mechanical statement about
molecular speeds, with no mention of temperature. But experiment also gives $PV = Nk_BT$. Setting the
two expressions for the same $PV$ equal must reveal what temperature is in molecular terms. This is the
moment the molecular picture pays off.

\subsection{The kinetic interpretation of temperature}

\begin{derivbox}{What temperature measures}
\stepn{1} Rewrite the pressure result in terms of the mean translational kinetic energy per molecule:
\[
PV = \tfrac13 Nm\overline{v^2} = \tfrac23 N\cdot\underbrace{\tfrac12 m\overline{v^2}}_{\overline K}.
\]

\stepn{2} Compare with the experimental ideal-gas law $PV = Nk_BT$:
\[
\tfrac23 N\overline K = Nk_BT
\quad\Longrightarrow\quad
\boxed{\;\overline K = \tfrac12 m\overline{v^2} = \tfrac32 k_B T\;}
\]

\stepn{3} \textbf{Read it.} The average translational kinetic energy of a molecule is $\tfrac32k_BT$. It
depends on \emph{temperature alone}: not on pressure, not on volume, and, most strikingly, \textbf{not on
which gas}. At the same $T$, a heavy xenon atom and a light helium atom carry the same average kinetic
energy; the xenon simply moves more slowly.
\end{derivbox}

\why{Temperature is the average translational kinetic energy of the molecules, expressed in different
units: $T=\tfrac{2}{3k_B}\overline K$. At room temperature $\overline K=\tfrac32k_B(\SI{300}{K})
=\SI{6.2e-21}{J}$, about \SI{0.04}{eV}. Absolute zero is the temperature at which this classical
translational energy would vanish. And the claim that \chref{ch:thermo} had to take on trust, that the
internal energy of an ideal gas depends on its temperature alone, becomes a theorem: for a monatomic gas,
whose molecules have no energy other than translational kinetic energy, $U = N\overline K = \tfrac32 Nk_BT
= \tfrac32nRT$, with no $P$ or $V$ in sight.}

\begin{trapbox}{temperature belongs to the crowd, not to one molecule}
A single molecule has a kinetic energy, but not a temperature. Temperature is a property of the
\emph{average} over a great many molecules in equilibrium, and the molecules in a gas at one temperature
have a wide spread of energies (Section~\ref{sec:ktmaxwell}). Equally, the kinetic energy that
temperature measures is that of the random motion: a gas in a moving car is no hotter than the same gas
at rest, because the bulk motion of its centre of mass (CM) is not random and is not counted.
\end{trapbox}

The relation also shows why absolute zero cannot be passed and why classical kinetic theory cannot be the
last word there. Quantum mechanics shows that molecules keep a small residual motion even at $T=0$, so
``all motion stops'' is the classical idealisation, not the full truth. At the temperatures of this
book's gases, hundreds of kelvin, the classical picture is excellent.

\subsection{Avogadro's hypothesis, and where the number came from}

The kinetic meaning of temperature explains one of the oldest clues that matter is made of molecules.

\begin{defbox}{Avogadro's hypothesis (1811)}
\textbf{Equal volumes of all gases, at the same temperature and pressure, contain equal numbers of
molecules.} Avogadro proposed this with no idea \emph{how many} molecules there are, only that the
count is the same for every gas, light or heavy.
\end{defbox}

\begin{derivbox}{Why the count is the same for every gas}
\stepn{1} \textbf{Same temperature, same energy.} The average translational kinetic energy of a molecule,
$\tfrac12 m\overline{v^2} = \tfrac32 k_BT$, depends on $T$ \emph{alone}. Heavy molecules move slowly and
light ones fast, but at a common temperature every gas carries the same average energy per molecule.

\stepn{2} \textbf{Same pressure, same energy density.} From $P=\tfrac23 n_V\overline K$, two gases at the
same pressure and with the same $\overline K$ must have the same number density $n_V$.

\stepn{3} \textbf{Solve for the count.} Equivalently, $N = PV/k_BT$, and the right-hand side contains only
$P$, $V$, $T$ and the universal $k_B$, \emph{nothing about the particular gas}. Fixing $P$, $V$ and $T$
fixes $N$. Avogadro's hypothesis is now a theorem.
\end{derivbox}

\inwords{At a given temperature all gases share the same energy budget per molecule. A gas does not push
harder because its molecules are heavier: the heavier molecules move more slowly, and the two effects
cancel exactly in the pressure, leaving the molecule \emph{count} as the only thing a given $P$, $V$ and
$T$ can pin down.}

\begin{remarkbox}{How Avogadro's number was measured}
Avogadro never knew the value; it was pinned down almost a century later. Because $k_B = R/N_A$ and $R$
was already known from gas experiments, \emph{any} independent measurement of $k_B$ delivers
\[
N_A = \frac{R}{k_B} \approx \SI{6.022e23}{mol^{-1}}.
\]
The decisive measurement was \textbf{Brownian motion}, the random jitter of pollen-sized particles
suspended in water. Einstein (1905) predicted how far such a particle should wander as invisible
molecules kick it, with $k_B$ (and so $N_A$) in the formula; Perrin (1908) measured the wandering under a
microscope and read off $N_A$. Independent routes (electrolysis with the electron's charge, X-ray
crystallography, black-body radiation) converged on the same number, which finally convinced the last
doubters that atoms are real. Perrin received the 1926 Nobel Prize for this work.

Since 2019 (as \chref{ch:thermal} noted for the kelvin) neither constant is measured any more: the SI \emph{defines}
$N_A=\SI{6.02214076e23}{mol^{-1}}$ and $k_B=\SI{1.380649e-23}{J/K}$ exactly, and the mole and the kelvin
are built on them. Like the speed of light, they have been promoted from quantities we measure to
constants we fix. The value $k_B=\SI{1.38e-23}{J/K}$ used in this chapter is the exact value rounded.
\end{remarkbox}

\subsection{The root-mean-square speed}

The pressure and temperature results both involve $\overline{v^2}$, the average of the \emph{squared}
speed. Its square root is a natural measure of a typical molecular speed, the \textbf{root-mean-square
speed} $v_{\text{rms}}$; read the name backwards: \emph{square} each speed, take the \emph{mean} of the
squares, then take the \emph{root}. Squares are used because energy and pressure depend on $v^2$, and
because averaging a velocity component without squaring gives zero, the plus and minus values cancelling.

\begin{keybox}{The rms speed}
From $\tfrac12 m\overline{v^2} = \tfrac32k_BT$, using $M = N_A m$ and $R = N_A k_B$:
\[
\boxed{\;v_{\text{rms}} = \sqrt{\overline{v^2}} = \sqrt{\frac{3k_BT}{m}} = \sqrt{\frac{3RT}{M}}
= \sqrt{\frac{3P}{\rho}}\;}
\]
So $v_{\text{rms}} \propto \sqrt{T}$ and $v_{\text{rms}}\propto 1/\sqrt{M}$: hot means fast, light means
fast. Pair $m$ (kilograms per molecule) with $k_B$, or $M$ (kilograms per mole) with $R$; never mix them.
\end{keybox}

\begin{exbox}{The speeds of nitrogen and hydrogen}\label{ex:ktrms}
Find $v_{\text{rms}}$ for nitrogen ($M = \SI{28}{g/mol}$) and hydrogen ($M=\SI{2}{g/mol}$) at
\SI{300}{K}.

\stepn{1} Nitrogen, with $M$ converted to \si{kg/mol}:
$v_{\text{rms}} = \sqrt{\dfrac{3(8.314)(300)}{0.028}} = \sqrt{2.67\times10^5} \approx \SI{517}{m/s}$.

\stepn{2} Hydrogen: the same formula with $M=\SI{0.002}{kg/mol}$, or by scaling,
$v_{\text{rms}} = 517\sqrt{28/2} = 517\sqrt{14} \approx \SI{1934}{m/s}$.

\stepn{3} \textbf{Check against sound.} Half a kilometre per second is comparable to the speed of sound in
air at \SI{300}{K}, \SI{347}{m/s}, and it should be: sound is carried by these very molecules, so it
cannot outrun them. (The ratio is worked out as a problem type at the end of the chapter.)

\stepn{4} \textbf{Hydrogen and the atmosphere.} Even \SI{1.9}{km/s} is far below Earth's escape speed of
\SI{11.2}{km/s} (\chref{ch:gravitation}, Section~\ref{sec:escape}), so why has Earth lost its free
hydrogen? Two refinements are needed. Escape happens from the top of the atmosphere, where the gas is
about \SI{1000}{K} and hydrogen is mostly atomic ($M=\SI{1}{g/mol}$), giving
$v_{\text{rms}}=\sqrt{3(8.314)(1000)/0.001}\approx\SI{5.0}{km/s}$. And the speeds are spread out
(Section~\ref{sec:ktmaxwell}): a small fraction of atoms in the fast tail exceed \SI{11.2}{km/s} at any
moment, and over geological time that steady leak empties the atmosphere of hydrogen and helium. Nitrogen
at \SI{1000}{K} has $v_{\text{rms}}\approx\SI{0.94}{km/s}$, twelve times below escape, and its tail is far
too thin to leak. A widely used rule of thumb, stated here without proof, is that a planet keeps a gas
for billions of years only if the gas's $v_{\text{rms}}$ is below about one-sixth of the escape speed.
\ans{\SI{517}{m/s}; \SI{1934}{m/s}}
\end{exbox}

\begin{trapbox}{rms is not the mean, and neither is ``the'' speed}
Molecular speeds are spread over a broad distribution (Figure~\ref{fig:ktmaxwell}), so no single number
tells the whole story. Three summaries are in use, the most probable speed $v_p$, the mean speed $\bar v$
and $v_{\text{rms}}$, and they are \emph{not} equal. Pressure and energy involve $\overline{v^2}$, so
$v_{\text{rms}}$ appears in those formulas; but do not call it the average speed, because
$v_{\text{rms}} > \bar v > v_p$. The next section makes the distinction precise.
\end{trapbox}

% ================================================================
\section{The Maxwell--Boltzmann Distribution}\label{sec:ktmaxwell}

Not every molecule moves at $v_{\text{rms}}$. At any instant the molecules in a gas share out a wide range
of speeds, some nearly at rest, a few very fast, and collisions constantly reshuffle which molecule has
which speed. Yet the \emph{fraction} of molecules in each range of speed stays fixed, set by the
temperature. Maxwell found that fraction in 1860, and Boltzmann later showed why it takes the form it
does. To use it we need the language of distributions from \chref{ch:toolkit},
Section~\ref{sec:averages}: a distribution function $\phi(v)$ is defined so that $\phi(v)\,dv$ is the
fraction of molecules with speeds between $v$ and $v+dv$, the fractions add up to one,
$\int_0^\infty\phi(v)\,dv=1$, and the average of any quantity $q(v)$ is $\overline{q}=\int_0^\infty q(v)\,\phi(v)\,dv$. (The Foundations wrote a general distribution as $f(v)$ and an average as $\langle q\rangle$; in this chapter $f$ counts degrees of freedom, so the speed distribution is $\phi(v)$, and an average is marked by a bar.)

\begin{keybox}{The Maxwell--Boltzmann speed distribution}
In a gas of molecules of mass $m$ at temperature $T$, the fraction of molecules with speeds between $v$
and $v+dv$ is $\phi(v)\,dv$, with
\[
\boxed{\;\phi(v) = 4\pi\left(\frac{m}{2\pi k_BT}\right)^{\!3/2} v^2\,
\exp\!\left(-\frac{mv^2}{2k_BT}\right)\;}
\]
The \emph{number} of molecules in that range is $N\phi(v)\,dv$. The function $\phi$ has units of
inverse speed (\si{s/m}), so that $\phi\,dv$ is a pure number. Two factors compete over its shape:
\begin{itemize}
\item the $v^2$ factor \emph{grows} with speed: there are more ways to move fast than slowly (more
      directions, in the sense made precise below), so the curve rises from zero;
\item the \textbf{Boltzmann factor} $\exp(-mv^2/2k_BT)=\exp(-K/k_BT)$ \emph{decays} with speed: fast
      molecules carry more energy and are exponentially rarer, so the curve eventually falls.
\end{itemize}
Their product peaks at a single most probable speed and trails off in a long high-speed tail.
\end{keybox}

\begin{figure}[htbp]
\centering
\tikzfig{kinetic_maxwell}
\figcaption{fig:ktmaxwell}{The Maxwell--Boltzmann distribution for nitrogen, computed from the formula. At \SI{300}{K} the three summary speeds sit in the fixed order $v_p<\bar v<v_{\text{rms}}$, because the curve is lopsided: it starts at zero speed and has a long tail to the right. At \SI{900}{K} the peak moves to higher speed (by $\sqrt3$) and drops, keeping the area equal to one, and the fast tail beyond \SI{1000}{m/s} (shaded) grows from about one molecule in a hundred to nearly one in three.}
\end{figure}

\inwords{The rising $v^2$ is geometry (many directions in which to move quickly); the falling exponential
is energy economics (fast states are costly). Their product is where most molecules live. Heating the gas
weakens the exponential's grip, so the peak slides to higher speed and the tail fattens
(Figure~\ref{fig:ktmaxwell}).}

\subsection{Where the formula comes from}

\scopebox{The examinations require the shape of the distribution and the three characteristic speeds, not
the derivation of $\phi(v)$. The argument below shows where the formula comes from; one step in it is
stated without proof.}

\begin{derivbox}{Maxwell's argument}
\stepn{1} \textbf{One component.} Let $g(v_x)\,dv_x$ be the fraction of molecules whose $x$-component of
velocity lies between $v_x$ and $v_x+dv_x$. Boltzmann's result, which belongs to statistical mechanics
and is \emph{stated here without proof}, is that in equilibrium the probability of a state of energy $E$
is proportional to $e^{-E/k_BT}$. For the $x$-motion, $E=\tfrac12mv_x^2$, so
\[
g(v_x) = \sqrt{\frac{m}{2\pi k_BT}}\;\exp\!\left(-\frac{m v_x^2}{2k_BT}\right),
\]
a bell-shaped (Gaussian) curve centred on $v_x=0$. The prefactor makes $\int_{-\infty}^{\infty}
g\,dv_x=1$, using the Gaussian integral $\int_{-\infty}^{\infty}e^{-\alpha u^2}du=\sqrt{\pi/\alpha}$
with $\alpha=m/2k_BT$. (Maxwell's own route assumed instead that the three components are independent and
that the distribution depends only on the speed; only an exponential of $v_x^2$ satisfies both.)

\stepn{2} \textbf{Three components.} The components are independent, so the fraction with velocity in a
small box $dv_x\,dv_y\,dv_z$ is the product
$g(v_x)g(v_y)g(v_z)\,dv_x\,dv_y\,dv_z
=\bigl(\tfrac{m}{2\pi k_BT}\bigr)^{3/2}e^{-mv^2/2k_BT}\,dv_x\,dv_y\,dv_z$,
which depends only on the speed $v$.

\stepn{3} \textbf{From velocity to speed.} Picture each velocity as a point in a space whose axes are
$v_x$, $v_y$, $v_z$. All velocities with speed between $v$ and $v+dv$ fill a thin spherical shell of
radius $v$ and thickness $dv$, of volume $4\pi v^2\,dv$. Multiplying the density of Step~2 by this volume
gives $\phi(v)\,dv$, and the formula in the box follows. The $4\pi v^2$ is the ``more ways to move fast''
factor: a larger shell holds more velocity directions.
\end{derivbox}

\subsection{The three characteristic speeds}

Each characteristic speed is a different question asked of the same distribution. The integrals needed are
the standard Gaussian moments of \chref{ch:toolkit}, Section~\ref{sec:averages}: with $a>0$,
\[
\int_0^\infty v^2e^{-av^2}dv=\frac{\sqrt\pi}{4}a^{-3/2},\qquad
\int_0^\infty v^3e^{-av^2}dv=\frac{1}{2a^2},\qquad
\int_0^\infty v^4e^{-av^2}dv=\frac{3\sqrt\pi}{8}a^{-5/2}.
\]

\begin{derivbox}{The three speeds from $\phi(v)$}
Write $a=m/2k_BT$, so that $\phi(v)=4\pi(a/\pi)^{3/2}v^2e^{-av^2}$. (The first integral above confirms
that $\int_0^\infty\phi\,dv=4\pi(a/\pi)^{3/2}\cdot\frac{\sqrt\pi}{4}a^{-3/2}=1$.)

\stepn{1} \textbf{Most probable speed} $v_p$: where is the peak? Set $d\phi/dv=0$. Only $v^2e^{-av^2}$
varies, and its derivative is $(2v-2av^3)e^{-av^2}$, which vanishes when $v^2=1/a$:
\[ v_p = \sqrt{\frac{2k_BT}{m}} = \sqrt{\frac{2RT}{M}}. \]

\stepn{2} \textbf{Mean speed} $\bar v$: what is the ordinary average?
\[
\bar v = \int_0^\infty v\,\phi(v)\,dv = 4\pi\Bigl(\frac{a}{\pi}\Bigr)^{3/2}\frac{1}{2a^2}
=\frac{2}{\sqrt{\pi a}}
=\sqrt{\frac{8k_BT}{\pi m}} = \sqrt{\frac{8RT}{\pi M}}.
\]

\stepn{3} \textbf{Root-mean-square speed} $v_{\text{rms}}$: what is $\sqrt{\overline{v^2}}$?
\[
\overline{v^2} = \int_0^\infty v^2 \phi(v)\,dv = 4\pi\Bigl(\frac{a}{\pi}\Bigr)^{3/2}\frac{3\sqrt\pi}{8}a^{-5/2}
=\frac{3}{2a}=\frac{3k_BT}{m},
\]
the same $\overline{v^2}$ that the pressure derivation and the kinetic meaning of temperature required.
So $v_{\text{rms}} = \sqrt{3k_BT/m} = \sqrt{3RT/M}$, and the distribution is consistent with
Section~\ref{sec:kttemperature}. Three questions (the peak, the average, the square root of the average
square) give three different numbers.
\end{derivbox}

\begin{keybox}{The three characteristic speeds}
All three scale as $\sqrt{RT/M}$ and differ only by a numerical factor:
\[
v_p = \sqrt{\frac{2RT}{M}}, \qquad
\bar v = \sqrt{\frac{8RT}{\pi M}}, \qquad
v_{\text{rms}} = \sqrt{\frac{3RT}{M}},
\]
in the ratio $\sqrt2 : \sqrt{8/\pi} : \sqrt3 = 1 : 1.128 : 1.225$, the same for every gas at every
temperature. Always $v_p < \bar v < v_{\text{rms}}$.
\begin{itemize}
\item $v_p$, the \textbf{most probable} speed: the peak of the curve, the commonest speed.
\item $\bar v$, the \textbf{mean} speed: the ordinary average, used for collision rates, effusion and the
      mean free path.
\item $v_{\text{rms}}$: the speed for energy and pressure, since those depend on $\overline{v^2}$.
\end{itemize}
\end{keybox}

\inwords{The curve is lopsided because speed cannot be negative (it is bounded below by zero) but has no
ceiling, so the tail can only stretch to the right. That asymmetry is exactly why the three averages
differ: for a symmetric bell curve the peak and the mean would coincide.}

\why{The fast tail, small as it is, runs much of the world. Only molecules in the high-speed tail have
enough energy to react, so raising the temperature, which stretches the tail dramatically, speeds up
chemical reactions and cooking. The same tail lets the fastest water molecules escape a puddle far below
boiling point (evaporation, which is why it cools what it leaves behind), and lets the fastest atmospheric
molecules exceed escape speed. Figure~\ref{fig:ktmaxwell} shows the scale of the effect: tripling the
temperature of nitrogen multiplies the fraction faster than \SI{1000}{m/s} by nearly thirty.}

\subsection{Effusion: leaking through a pinhole}

A gas leaking through a hole much smaller than the mean free path (Section~\ref{sec:ktmfp}) does not flow
as a fluid; molecules simply pass through one at a time when their random motion carries them into the
hole. This is \textbf{effusion}. The number of molecules crossing unit area of the hole per second is
\[
\Phi = \tfrac14\,n_V\,\bar v ,
\]
stated here without proof: the factor $\tfrac14$ comes from averaging over the directions of approach,
which needs an integral over angles beyond this chapter's scope. Its consequence is easy to read: at the
same temperature and pressure (so the same $n_V$), the effusion rate is proportional to $\bar v$, that is,
to $1/\sqrt{M}$. This is \textbf{Graham's law},
\[
\frac{\text{rate}_A}{\text{rate}_B}=\sqrt{\frac{M_B}{M_A}} ,
\]
found experimentally by Graham in 1846 and explained by kinetic theory. The lighter gas escapes faster,
which is how the isotopes of uranium were first separated on an industrial scale (a challenge problem at
the end of the chapter works out the numbers).

% ================================================================
\section{Equipartition and Specific Heats}\label{sec:ktequipartition}

So far every molecule has been treated as a structureless point that can only fly in straight lines.
For a monatomic gas such as helium that is the whole story, giving an average energy
$\overline K = \tfrac32 k_BT$ per molecule and $U=\tfrac32nRT$ per sample. But a nitrogen molecule is a dumbbell: it can also \emph{tumble}, and tumbling stores
energy too. To predict the heat capacities of real gases we must count \emph{all} the ways a molecule can
hold energy, not just its three directions of motion, and find how energy is shared among them.

\subsection{Degrees of freedom}

\begin{defbox}{Degrees of freedom $f$}
For equipartition, the number of \textbf{degrees of freedom} $f$ of a molecule is the number of
independent \emph{squared} terms in its energy. Each direction of translation contributes a
$\tfrac12mv_x^2$-type term; each axis of rotation contributes a $\tfrac12I\omega^2$-type term, where $I$
is the moment of inertia about that axis and $\omega$ the angular speed about it (\chref{ch:rotation}).
\begin{itemize}
\item \textbf{Monatomic} (He, Ar): three translations only: $f = 3$.
\item \textbf{Diatomic, rigid} (N$_2$, O$_2$ at ordinary temperatures): 3 translations $+$ 2 rotations
      about the two axes perpendicular to the bond: $f = 5$.
\item \textbf{Linear polyatomic, rigid} (CO$_2$): also $f=5$, for the same reason.
\item \textbf{Non-linear polyatomic, rigid} (CH$_4$, H$_2$O vapour): 3 translations $+$ 3 rotations:
      $f = 6$.
\end{itemize}
A \textbf{vibration} (the atoms oscillating along a bond like masses on a spring) contributes
\emph{two} squared terms, a kinetic $\tfrac12\mu\dot x^2$ and a potential $\tfrac12kx^2$ (with $\mu$ an
effective mass, $x$ the stretch of the bond and $k$ the bond's spring constant, not to be confused with
$k_B$). A diatomic molecule that vibrates therefore has $f=7$.
\end{defbox}

\begin{trapbox}{why some motions do not count}
The usual explanation for leaving out rotation about the bond axis of N$_2$, that its moment of inertia is
almost zero, is incomplete. Classically, the average energy in a squared term does not depend on the size
of $I$ at all (the derivation below shows $m$ or $I$ cancelling), so a tiny $I$ would still earn its share.
The real reason is quantum mechanical: rotational and vibrational energy can only be taken up in steps,
and when the smallest step is much larger than $k_BT$ the motion cannot be excited and is ``frozen out''.
For rotation about the bond the step is enormous because $I$ is tiny; for the vibration of N$_2$ the step
corresponds to a temperature of about \SI{3400}{K}. So $f$ is not a fixed property of a molecule: it
depends on the temperature.
\end{trapbox}

\subsection{The law of equipartition of energy}

We found that the translational energy of a molecule is $\tfrac32 k_BT$. By isotropy it is split equally,
$\tfrac12 k_BT$ along each of $x$, $y$ and $z$. The law of equipartition says this even-handedness is
universal: \emph{every} squared term gets the same $\tfrac12 k_BT$ on average. Here is why one
translational direction gets exactly that amount, by an argument that applies unchanged to rotation and
vibration.

\begin{derivbox}{Why one degree of freedom carries $\tfrac12 k_BT$}
\stepn{1} \textbf{The distribution of one velocity component} is the Gaussian $g(v_x)$ of Maxwell's
argument (Section~\ref{sec:ktmaxwell}), whose exponent is minus the kinetic energy in $x$ divided by
$k_BT$.

\stepn{2} \textbf{The average $x$-kinetic energy} is that energy weighted by how often each $v_x$
occurs:
\[
\overline{\tfrac12 m v_x^2}
= \int_{-\infty}^{\infty} \tfrac12 m v_x^2\; g(v_x)\,dv_x .
\]

\stepn{3} \textbf{Evaluate the Gaussian integral.} With
$\int_{-\infty}^{\infty} u^2\,e^{-\alpha u^2}\,du = \tfrac12\sqrt{\pi}\,\alpha^{-3/2}$ and
$\alpha = m/2k_BT$,
\[
\overline{\tfrac12 m v_x^2}=\tfrac12 m\sqrt{\frac{\alpha}{\pi}}\cdot\tfrac12\sqrt\pi\,\alpha^{-3/2}
=\frac{m}{4\alpha}=
\boxed{\;\tfrac12 k_BT\;}
\]
The mass has cancelled: heavy and light molecules carry the \emph{same} average energy per direction,
which is the result of Section~\ref{sec:kttemperature}, now seen to hold one axis at a time.

\stepn{4} \textbf{Generalise.} Nothing in Steps 1--3 used the fact that the energy was
\emph{translational}. Any energy that is \textbf{quadratic} in its variable (a rotational $\tfrac12
I\omega^2$, or a vibrational $\tfrac12 k x^2$) carries the same Gaussian Boltzmann factor and therefore
the same average, $\tfrac12 k_BT$, with $I$ or $k$ cancelling as $m$ did. That is the law of
equipartition.
\end{derivbox}

\begin{keybox}{The law of equipartition of energy}
In thermal equilibrium at temperature $T$, each independent \textbf{quadratic} term in a molecule's
energy carries, on average,
\[
\boxed{\;\tfrac12 k_BT \text{ per molecule}\quad\bigl(=\tfrac12 RT \text{ per mole}\bigr)\;}
\]
So a molecule with $f$ degrees of freedom has average energy $\tfrac{f}{2}k_BT$. Check: a monatomic gas
has $f=3$ and $\tfrac32 k_BT$, the result of Section~\ref{sec:kttemperature}. The law is classical: it
holds only for motions whose quantum steps are small compared with $k_BT$.
\end{keybox}

\why{Equipartition is why the heat capacity of a gas betrays its molecular shape. A monatomic gas can only
\emph{move}; a diatomic one can also \emph{tumble}; each new way to store energy adds another $\tfrac12 R$
per mole to the heat capacity. So measuring how much heat a gas needs to warm by one kelvin tells you how
many atoms are in each molecule and roughly how they are arranged: thermodynamics reading out the
invisible architecture of matter.}

\subsection{Predicting the specific heats}

\begin{derivbox}{$C_v$, $C_p$ and $\gamma$ from molecular structure}
\stepn{1} The internal energy of one mole with $f$ degrees of freedom (an ideal gas has no potential
energy between molecules, assumption (iii)) is
\[
U = N_A \cdot f\cdot\tfrac12 k_BT = \tfrac{f}{2}RT .
\]

\stepn{2} The molar heat capacity at constant volume is $C_v = dU/dT$ (\chref{ch:thermo}: at constant
volume no work is done, so all the heat goes into $U$):
\[
C_v = \frac{f}{2}R .
\]

\stepn{3} Mayer's relation $C_p=C_v+R$ (\chref{ch:thermo}) then fixes the rest, with $\gamma=C_p/C_v$
the ratio of the molar heat capacities:
\[
C_p = \left(\frac{f}{2}+1\right)R,
\qquad
\gamma = \frac{C_p}{C_v} = 1 + \frac{2}{f}.
\]
\end{derivbox}

\begin{tablebox}{Predicted and measured molar heat capacities (\si{J.mol^{-1}.K^{-1}})}
\centering
\begin{tabular}{@{}lccccc@{}}
\toprule
 & $f$ & $C_v$ predicted & $C_p$ predicted & $\gamma$ predicted & $C_v$ measured (near \SI{300}{K})\\
\midrule
Monatomic & 3 & $\tfrac32R = 12.5$ & $\tfrac52R=20.8$ & $5/3 \approx 1.67$ & He 12.5, Ar 12.5\\
Diatomic (rigid) & 5 & $\tfrac52R = 20.8$ & $\tfrac72R=29.1$ & $7/5 = 1.40$ & N$_2$ 20.8, O$_2$ 21.1, H$_2$ 20.5\\
Non-linear (rigid) & 6 & $3R=24.9$ & $4R=33.3$ & $4/3 \approx 1.33$ & H$_2$O vapour about 25\\
Linear CO$_2$ (rigid) & 5 & $\tfrac52R = 20.8$ & $\tfrac72R=29.1$ & $1.40$ & CO$_2$ 28.8\\
\bottomrule
\end{tabular}

\smallskip\raggedright
The monatomic and simple diatomic gases agree to within about \SI{2}{\percent}. Carbon dioxide does not:
its bending vibrations need only a small quantum step and are already partly excited at room temperature,
so it stores more energy per kelvin than the rigid model allows.
\end{tablebox}

\inwords{More degrees of freedom means more ``pockets'' into which heat can be tucked without raising the
translational energy that a thermometer responds to. So complex molecules need more heat per kelvin, and
their $\gamma$ lies closer to 1.}

\begin{remarkbox}{Hydrogen's heat capacity changes with temperature}
The measured $C_v$ of hydrogen gas is only about $\tfrac32R$ at a few tens of kelvin, rises to $\tfrac52R$ near
room temperature, and climbs towards $\tfrac72R$ at a few thousand kelvin. Rotation switches on as $k_BT$
reaches the rotational step, vibration only at far higher temperature. Classical equipartition, which
gives one fixed $f$, cannot explain the steps at all; Maxwell called this the greatest difficulty facing
the molecular theory. Quantum mechanics explained it in the twentieth century.
\end{remarkbox}

\subsection{From molar heat capacity to specific heat}

The derivation gives \emph{molar} heat capacities, in joules per \emph{mole} per kelvin. But engineers
usually work with mass, and want the \textbf{specific heat capacity} $s$ of \chref{ch:thermal}: heat per
\emph{kilogram} per kelvin. The two are one short step apart.

\begin{keybox}{Molar heat capacity $C$ and specific heat $s$}
A mole has mass $M$ (in \si{kg/mol}), so heating one mole is the same as heating $M$ kilograms:
\[
\boxed{\;C = M s \qquad\Longleftrightarrow\qquad s = \frac{C}{M}\;}
\]
Because $C$ is nearly the same for all gases of a given molecular shape, the \emph{specific} heat
$s = C/M$ is large for light gases and small for heavy ones: hydrogen stores far more heat per kilogram
than argon.
\end{keybox}

\begin{exbox}{The specific heat of air}\label{ex:ktair}
Predict the specific heat of air at constant pressure. Treat air as a rigid diatomic gas ($f=5$) with
average molar mass $M = \SI{28.8}{g/mol}$.

\stepn{1} Molar heat capacity: $C_p = \tfrac72 R = 3.5(8.314) = \SI{29.1}{J.mol^{-1}.K^{-1}}$.

\stepn{2} Divide by the molar mass to get the value per kilogram:
\[
s_p = \frac{C_p}{M} = \frac{29.1}{0.0288} \approx \SI{1010}{J.kg^{-1}.K^{-1}}.
\]

\stepn{3} The measured value is about \SI{1005}{J.kg^{-1}.K^{-1}}, agreement to about half a per cent from
nothing but counting a diatomic molecule's five degrees of freedom. It is a quarter of water's
\SI{4186}{J.kg^{-1}.K^{-1}}, and air is about 800 times less dense than water, so a cubic metre of air
stores some three thousand times less heat per kelvin than a cubic metre of water. That is why heating a
room's air is cheap, and why it is the sea, not the air, that moderates a coastal climate.
\ans{$s_p \approx \SI{1010}{J.kg^{-1}.K^{-1}}$}
\end{exbox}

\begin{remarkbox}{Solids and the law of Dulong and Petit}
The same counting works for a solid. An atom in a crystal oscillates about its lattice site, held by its
neighbours as if by springs (\chref{ch:solids}): three kinetic and three potential squared terms, so
$f=6$ and $U = 3RT$ per mole. That gives $C \approx 3R \approx \SI{25}{J.mol^{-1}.K^{-1}}$ for
\emph{nearly every} elemental solid at room temperature, the empirical law of Dulong and Petit (1819), and
explains the specific heats tabulated in \chref{ch:thermal}: for copper ($M=\SI{63.5}{g/mol}$),
$s=24.9/0.0635\approx\SI{392}{J.kg^{-1}.K^{-1}}$, against a measured \SI{385}{J.kg^{-1}.K^{-1}}. The law
fails at low temperatures, and for light, stiff solids even at room temperature (diamond has
$C\approx\SI{6}{J.mol^{-1}.K^{-1}}$), for the quantum reason given above. That failure was one of the
first experimental cracks through which quantum theory entered physics.
\end{remarkbox}

% ================================================================
\section{Mean Free Path}\label{sec:ktmfp}

If molecules fly at \SI{500}{m/s}, why does a smell take so long to cross a still room? Because no
molecule flies far in a straight line: it collides, changes direction, and staggers along a random path.
The \textbf{mean free path} $\lambda$ is the average distance a molecule travels between successive
collisions. Estimating it needs one more property of a molecule: its size. Model each molecule as a hard
sphere of diameter $d$.

\begin{figure}[htbp]
\centering
\tikzfig{kinetic_cylinder}
\figcaption{fig:ktcylinder}{The collision cylinder. Two spheres of diameter $d$ touch when their centres are $d$ apart, so the moving molecule hits every molecule whose \emph{centre} lies within $d$ of its path: a cylinder of radius $d$ and cross-section $\pi d^2$, four times the molecule's own cross-section. Counting the centres inside the cylinder counts the collisions.}
\end{figure}

\begin{derivbox}{The mean free path}
\stepn{1} \textbf{The target area.} A molecule hits any other whose centre comes within a distance $d$ of
its own centre (Figure~\ref{fig:ktcylinder}). So it sweeps out a cylinder of radius $d$, whose cross-section $\pi d^2$ is called the \textbf{collision cross-section}.

\stepn{2} \textbf{First estimate: targets at rest.} Moving at speed $\bar v$ for a time $t$, the molecule
sweeps a cylinder of volume $\pi d^2\bar v t$. With $n_V$ molecules per unit volume, it meets
$n_V\pi d^2\bar v t$ of them. (After each collision it changes direction, so the cylinder is really
bent at each hit, but its volume is the same.) The distance travelled per collision is
\[
\lambda \approx \frac{\bar v t}{n_V\pi d^2 \bar v t} = \frac{1}{\pi d^2 n_V}.
\]

\stepn{3} \textbf{The targets move too.} How often collisions happen depends on how fast molecules
approach each other, the \emph{relative} speed $|\vv_1-\vv_2|$. Its square is
$v_1^2+v_2^2-2\vv_1\cdot\vv_2$, and because the directions are random, $\vv_1\cdot\vv_2$ averages to zero.
So the mean-square relative speed is $2\overline{v^2}$, and the typical relative speed is $\sqrt2$ times
the typical speed. For the Maxwell distribution the mean relative speed is exactly $\sqrt2\,\bar v$
(stated without proof: the relative velocity of two molecules is itself Maxwell-distributed, with half
the mass). The number of collisions in time $t$ is therefore $n_V\pi d^2(\sqrt2\,\bar v)t$, while the
distance actually travelled is still $\bar vt$:
\[
\boxed{\;\lambda = \frac{1}{\sqrt{2}\,\pi d^2 n_V}\;}
\]
This book always uses this form, with the $\sqrt2$. Some books quote $1/\pi d^2n_V$, which is the estimate
with targets at rest; it is about \SI{40}{\percent} too large.
\end{derivbox}

\inwords{$\lambda$ falls as the gas gets denser ($\propto 1/n_V$) and as the molecules get bigger
($\propto 1/d^2$). At a fixed temperature $n_V \propto P$, so halving the pressure doubles the free path.
That is why vacuum equipment works: pump the pressure down far enough and molecules fly from wall to wall
without meeting each other at all.}

\begin{keybox}{The mean free path in terms of pressure, and the collision frequency}
The number density is awkward to measure directly, but $P$ and $T$ are not. Substituting
$n_V = P/k_BT$ turns the free path into laboratory variables:
\[
\boxed{\;\lambda = \frac{k_BT}{\sqrt2\,\pi d^2 P}\;}
\]
so $\lambda \propto T$ at fixed pressure and $\lambda \propto 1/P$ at fixed temperature. The
\textbf{collision frequency} $z$, the number of collisions each molecule suffers per second, is the
distance travelled per second divided by the distance per collision:
\[
\boxed{\;z = \frac{\bar v}{\lambda} = \sqrt2\,\pi d^2 n_V\bar v\;}
\]
\end{keybox}

\begin{trapbox}{the mean free path does not depend on speed}
$\lambda=1/\sqrt2\pi d^2n_V$ contains no $\bar v$: a faster molecule meets molecules more often, but
covers ground faster in exact proportion. Heating a \emph{sealed} rigid container (fixed $n_V$) leaves
$\lambda$ unchanged; the molecules simply run their unchanged free paths more quickly, so $z$ rises as
$\sqrt T$. Only at fixed \emph{pressure}, when heating also thins the gas, does $\lambda$ grow with $T$.
\end{trapbox}

\begin{exbox}{Free path and collision rate in air}\label{ex:ktmfp}
Estimate $\lambda$ and $z$ for air at \SI{0}{\celsius} and atmospheric pressure, taking
$d = \SI{3.7e-10}{m}$ (a value found from the viscosity of air), $n_V = \SI{2.69e25}{m^{-3}}$, and
$M=\SI{29}{g/mol}$.

\stepn{1} Cross-section: $\pi d^2 = \pi(3.7\times10^{-10})^2 = \SI{4.30e-19}{m^2}$.

\stepn{2} Mean free path:
\[
\lambda = \frac{1}{\sqrt2 \times (4.30\times10^{-19}) \times (2.69\times10^{25})}
= \frac{1}{1.63\times10^{7}} \approx \SI{6.1e-8}{m} = \SI{61}{nm}.
\]

\stepn{3} Mean speed: $\bar v=\sqrt{8RT/\pi M}=\sqrt{8(8.314)(273)/(\pi\times0.029)}\approx\SI{447}{m/s}$, so
$z=\bar v/\lambda\approx 447/(6.1\times10^{-8})\approx\SI{7e9}{s^{-1}}$.

\stepn{4} \textbf{Perspective.} A molecule travels about $\lambda/d\approx165$ of its own diameters between
hits, passing some twenty neighbours (spaced \SI{3.3}{nm} apart) before striking one, and it collides
about seven billion times a second. That collision rate is why a gas settles into equilibrium, and into
the Maxwell distribution, almost instantly.
\ans{$\lambda\approx\SI{61}{nm}$; $z\approx\SI{7e9}{s^{-1}}$}
\end{exbox}

\begin{remarkbox}{Why a hot gas is more viscous}
\chref{ch:fluids} noted that a gas, unlike a liquid, grows \emph{more} viscous when heated. The mean free
path explains it. Viscosity is the drag of one layer of fluid on its neighbour. In a gas the layers do not
pull on each other; instead molecules fly across from one layer to the next, about a free path $\lambda$
at a time, carrying the flow velocity of the layer they left. A rough version of that count (stated here
without the full calculation) gives a coefficient of viscosity of about
$\tfrac13\rho\bar v\lambda$, where $\rho$ is the gas's density; a careful treatment of hard spheres raises
the $\tfrac13$ to about $\tfrac12$. Two predictions follow at once. Since $\rho\propto n_V$ and
$\lambda\propto1/n_V$, their product does not depend on the density, so the viscosity of a gas is
\emph{independent of its pressure}, a result so surprising that Maxwell checked it experimentally himself.
And since $\bar v\propto\sqrt T$, the viscosity rises with temperature, the opposite of a liquid, whose
viscosity comes from molecules held by their neighbours' attraction. For air at \SI{0}{\celsius},
$\tfrac12\rho\bar v\lambda=\tfrac12(1.29)(447)(\num{6.1e-8})\approx\SI{1.8e-5}{Pa.s}$, close to the measured \SI{1.7e-5}{Pa.s}; run backwards, this is how the molecular diameter in Example~\ref{ex:ktmfp} was found.
\end{remarkbox}

\begin{remarkbox}{Why smells cross a room by draughts, not by diffusion}
A molecule that takes random steps of length $\lambda$ in random directions does not get far. After $N$
steps its typical distance from the start grows only as $\lambda\sqrt N$, not $N\lambda$ (stated without
proof; it is the random-walk result). In a time $t$ it takes $N=z t$ steps, so it wanders a distance of
about $\lambda\sqrt{zt}=\sqrt{\lambda\bar v\,t}$. For air,
$\sqrt{(6.1\times10^{-8})(447)(\SI{1}{s})}\approx\SI{5}{mm}$ in the first second, and crossing a
\SI{5}{m} room this way would take $t\approx(\SI{5}{m})^2/\lambda\bar v\approx 9\times10^5$~s, about ten
days (an order-of-magnitude estimate). A smell actually crosses a room in a minute or so because the air
itself moves, in draughts and convection currents. Diffusion matters on small scales: across a cell, or
through the thin layer of still air at a leaf's surface.
\end{remarkbox}

% ================================================================
\section{Summary}\label{sec:ktsummary}
Kinetic theory explains the gas laws as Newtonian mechanics applied to a vast number of molecules.
Pressure is momentum delivered to the walls per unit time and area; comparing it with the ideal-gas law
shows that temperature measures the average translational kinetic energy of a molecule; equipartition
shares energy equally among the squared terms of a molecule's energy, which predicts the heat capacities;
and the size of the molecules sets how far they fly between collisions.

\begin{itemize}
\item Ideal gas: $PV=nRT=Nk_BT$, $P=n_Vk_BT$, with $k_B=R/N_A=\SI{1.38e-23}{J/K}$ and
      $n_V=N/V$.
\item Pressure from impacts: $P=\tfrac13\rho\overline{v^2}=\tfrac23n_V\overline K$; in a mixture the
      partial pressures add (Dalton).
\item Temperature: $\overline K=\tfrac12m\overline{v^2}=\tfrac32k_BT$, the same for every gas at the same
      $T$; so equal volumes at equal $P$ and $T$ hold equal numbers of molecules (Avogadro).
\item Maxwell--Boltzmann distribution $\phi(v)\propto v^2e^{-mv^2/2k_BT}$, with
      $v_p=\sqrt{2RT/M}<\bar v=\sqrt{8RT/\pi M}<v_{\text{rms}}=\sqrt{3RT/M}$, in the ratio
      $1:1.128:1.225$.
\item Effusion rate $\propto\bar v\propto1/\sqrt M$ at fixed $T$ and $P$ (Graham).
\item Equipartition: $\tfrac12k_BT$ per squared term; $U=\tfrac f2nRT$, $C_v=\tfrac f2R$,
      $C_p=C_v+R$, $\gamma=1+\tfrac2f$, with $f=3$, 5, 6 for monatomic, rigid linear and rigid non-linear
      molecules (each active vibration adds 2). Specific heat $s=C/M$; solids $C\approx3R$.
\item Mean free path $\lambda=\dfrac{1}{\sqrt2\pi d^2n_V}=\dfrac{k_BT}{\sqrt2\pi d^2P}$; collision
      frequency $z=\bar v/\lambda$.
\item Limits: dense or cold gases (finite molecular size, intermolecular forces), and low temperatures,
      where quantum steps freeze out rotation and vibration and equipartition fails.
\end{itemize}

% ================================================================
\section{Problem Types}\label{sec:kttypes}
The kinetic-theory questions examiners set fall into eleven types, six of them on molecular speeds. Almost
all reduce to a single proportionality, so the most efficient approach is to take ratios and let the
constants cancel. Because this chapter closes the volume, the habits that decide problems throughout
Volume~II are collected first.

\begin{methodbox}{Four habits for every problem in this volume}
\begin{enumerate}
\item \textbf{Convert to kelvin} whenever a temperature enters a ratio or a product: rms speeds, gas
      laws, Carnot efficiencies and the $T^4$ of radiation all demand absolute temperature. Only a
      temperature \emph{difference} may stay in degrees Celsius.
\item \textbf{Decide what is held fixed}: which variable ($P$, $V$ or $T$), which phase, which process
      (isothermal or adiabatic, sealed or open). That one choice selects the formula.
\item \textbf{Use state functions as shortcuts}: $\Delta U$ depends only on the end states, so compute it
      on the easiest path; around a closed cycle it is zero, and the net work is the enclosed area.
\item \textbf{Watch for deliberate distractors}: an irrelevant pressure in an rms-speed question, an
      irrelevant length in a stress question, a mass ratio in an average-kinetic-energy question.
\end{enumerate}
Recognising the type from the first line of a problem tells you which habit to use, which is why naming
the pattern is half the solution. Practise by covering the Method line of each type below and
reconstructing it from the Tell alone.
\end{methodbox}

\archq{rms speed scaling}
{``A gas has rms speed $v$ at $T_1$; find it at $T_2$,'' often with a \emph{pressure change} added.}
{$v_{\text{rms}}=\sqrt{3RT/M}$, so for the same gas $v_2=v_1\sqrt{T_2/T_1}$. Different gases at the same
$T$: $v_{\text{rms}}\propto1/\sqrt M$.}
{The pressure is a distractor: the rms speed depends only on $T$ and $M$. Temperatures must be in kelvin.}
{The rms speed of the molecules of a gas is \SI{200}{m/s} at \SI{27}{\celsius}. The gas is heated to
\SI{127}{\celsius} while its pressure is also changed. The new rms speed is:}
{\SI{173}{m/s}}
{\SI{267}{m/s}}
{\SI{231}{m/s}}
{\SI{434}{m/s}\archans{C}{$v_2=200\sqrt{400/300}=200\sqrt{4/3}\approx\SI{231}{m/s}$; the pressure change
is irrelevant. (A) inverts the temperature ratio; (B) takes $v\propto T$ instead of $\sqrt T$; (D) uses
the Celsius temperatures, $200\sqrt{127/27}$.}}

\archq{Average kinetic energy at the same temperature}
{``Two different gases at the same temperature: compare the average kinetic energy per molecule,'' often
with a mass ratio given.}
{The average translational kinetic energy per molecule is $\tfrac32 k_BT$, temperature alone. At the same
$T$, every gas has the same value.}
{The mass ratio is a pure distractor for kinetic energy per \emph{molecule}. (The \emph{total} energy per
molecule, including rotation, does depend on the atomicity.)}
{A vessel contains a monatomic gas and a diatomic gas in the mass ratio $4:1$ at \SI{300}{K}. The ratio of
the average translational kinetic energy \emph{per molecule} of the two gases is:}
{$1:1$}
{$4:1$}
{$3:5$}
{$1:4$\archans{A}{The average translational kinetic energy per molecule is $\tfrac32 k_BT$ for every
gas, so at the same $T$ the ratio is $1:1$. (B) and (D) use the mass ratio; (C) compares the total
energies $\tfrac32k_BT$ and $\tfrac52k_BT$, which is not what was asked.}}

\archq{The escape-speed temperature}
{``At what temperature would a molecule's rms speed equal the escape speed of a planet?''}
{Set $v_{\text{rms}}=v_{\text{esc}}$: $\tfrac12mv_{\text{esc}}^2=\tfrac32k_BT$, so
$T=\dfrac{m\,v_{\text{esc}}^2}{3k_B}$.}
{Pair the molecular mass $m$ with $k_B$, or the molar mass $M$ with $R$; never mix them. Use $3k_B$, from
$\tfrac32k_BT$, not $k_B$ or $2k_B$.}
{At what temperature would the rms speed of oxygen molecules (mass \SI{5.31e-26}{kg}) equal the escape
speed from Earth, \SI{11.2}{km/s}? (Take $k_B=\SI{1.38e-23}{J/K}$.)}
{\SI{8.0e4}{K}}
{\SI{2.4e5}{K}}
{\SI{4.8e5}{K}}
{\SI{1.6e5}{K}\archans{D}{$T=\dfrac{m\,v_{\text{esc}}^2}{3k_B}=\dfrac{(5.31\times10^{-26})(11200)^2}
{3(1.38\times10^{-23})}\approx\SI{1.6e5}{K}$. (A) uses $6k_B$, as if $\tfrac12mv^2$ were $3k_BT$; (B) uses $2k_B$, which
sets the \emph{most probable} speed equal to $v_{\text{esc}}$; (C) drops the factor 3. No planetary
atmosphere is this hot; oxygen is held because escape needs only a thin tail of the distribution, and for
oxygen even that tail is negligible.}}

\archq{Pressure when mass or speed is rescaled}
{``The molecular mass is halved and the speeds doubled; find the new pressure.''}
{From $P=\tfrac13n_Vm\overline{v^2}$, track each factor: $P\propto n_Vm\overline{v^2}$.}
{Speed enters \emph{squared}: doubling every speed multiplies $\overline{v^2}$ by four, not two.}
{A gas exerts pressure $P_0$ on its container. If the mass of each molecule is halved and the speed of each
is doubled, the number of molecules and the volume being unchanged, the new pressure is:}
{$P_0$}
{$4P_0$}
{$P_0/2$}
{$2P_0$\archans{D}{$P\propto m\overline{v^2}$, so the factor is $\tfrac12\times4=2$. (A) forgets that speed
enters squared ($\tfrac12\times2$); (B) ignores the halved mass; (C) applies only the mass change.}}

\archq{Effusion and Graham's law}
{``Two gases escape through a pinhole: compare their rates,'' or ``find an unknown molar mass from its
effusion rate.''}
{At the same $T$ and $P$, rate $\propto\bar v\propto1/\sqrt M$, so
$\dfrac{\text{rate}_A}{\text{rate}_B}=\sqrt{\dfrac{M_B}{M_A}}$.}
{The lighter gas is \emph{faster}: do not invert the ratio, and remember the square root.}
{An unknown gas effuses through a pinhole at \SI{2}{mL/min}, while hydrogen ($M=\SI{2}{g/mol}$) effuses
through the same hole at \SI{11}{mL/min} under identical conditions. The molar mass of the unknown gas is
about:}
{\SI{11}{g/mol}}
{\SI{60.5}{g/mol}}
{\SI{121}{g/mol}}
{\SI{0.066}{g/mol}\archans{B}{$\dfrac{2}{11}=\sqrt{\dfrac{2}{M}}$, so
$M=2\left(\dfrac{11}{2}\right)^2=\SI{60.5}{g/mol}$. (A) drops the square root (rate $\propto1/M$); (C)
takes hydrogen's molar mass as 4; (D) inverts the ratio, making the slower gas the lighter.}}

\archq{The speed of sound and the rms speed}
{``Relate the speed of sound in a gas to the rms speed of its molecules.''}
{The speed of sound in a gas is $\sqrt{\gamma P/\rho}=\sqrt{\gamma RT/M}$ (quoted here: waves are not
treated in this volume; the $\gamma$ appears because the compressions of a sound wave are too quick for
heat to flow, so they are adiabatic). With $v_{\text{rms}}=\sqrt{3RT/M}$,
$\dfrac{v_{\text{sound}}}{v_{\text{rms}}}=\sqrt{\dfrac{\gamma}{3}}$, depending on neither $T$ nor $P$.}
{Use the adiabatic $\gamma P$, not the isothermal $P$ (Newton's historical error).}
{For air ($\gamma=1.4$), the ratio of the speed of sound to the rms speed of the molecules is closest to:}
{$0.68$}
{$0.58$}
{$1.00$}
{$1.46$\archans{A}{$\sqrt{\gamma/3}=\sqrt{1.4/3}\approx0.68$, so sound is somewhat slower than the rms
molecular speed, as it must be, since the molecules carry it. (B) is Newton's isothermal value
$\sqrt{1/3}$; (C) assumes the two are equal; (D) is the ratio inverted.}}

\archq{The three speeds of the Maxwell distribution}
{``Given one of $v_p$, $\bar v$, $v_{\text{rms}}$, find another,'' or ``put them in order.''}
{All three are $\sqrt{RT/M}$ times $\sqrt2$, $\sqrt{8/\pi}$ and $\sqrt3$, in the ratio
$1:1.128:1.225$, so convert by ratios and never compute $RT/M$.}
{Mixing up which is largest: the order is always $v_p<\bar v<v_{\text{rms}}$.}
{The rms speed of the molecules of a gas is \SI{500}{m/s}. Their most probable speed is:}
{\SI{500}{m/s}}
{\SI{461}{m/s}}
{\SI{408}{m/s}}
{\SI{612}{m/s}\archans{C}{$v_p=v_{\text{rms}}\sqrt{2/3}=500\times0.816\approx\SI{408}{m/s}$. (B) is the
mean speed, $500\sqrt{8/3\pi}$; (D) multiplies by $\sqrt{3/2}$ instead of dividing; (A) treats the three
speeds as equal.}}

\archq{Internal energy of a gas mixture}
{``A mixture of $n_1$ moles of one gas and $n_2$ of another (different atomicities) at $T$; find the
total internal energy.''}
{Each gas contributes $n_i\dfrac{f_i}{2}RT$; add them. Monatomic $f=3$, rigid diatomic $f=5$.}
{Use each gas's \emph{own} degrees of freedom; do not average them first. ``Neglect vibration'' means the
diatomic molecules are rigid ($f=5$).}
{\SI{2}{mol} of oxygen (diatomic) is mixed with \SI{4}{mol} of argon (monatomic) at temperature $T$.
Neglecting vibration, the internal energy of the mixture is:}
{$9RT$}
{$12RT$}
{$11RT$}
{$13RT$\archans{C}{$U=2\left(\tfrac52 RT\right)+4\left(\tfrac32 RT\right)=5RT+6RT=11RT$. (A) counts only
translation for both; (B) averages $f$ to 4 before multiplying by 6 moles; (D) swaps the two values of
$f$.}}

\archq{$C_v$ and $\gamma$ of a mixture}
{``Mix $n_1$ moles of one gas with $n_2$ of another; find $C_v$, $C_p$ or $\gamma$ of the mixture.''}
{Heat capacities add: $C_v^{\text{mix}}=\dfrac{n_1C_{v1}+n_2C_{v2}}{n_1+n_2}$, then
$C_p^{\text{mix}}=C_v^{\text{mix}}+R$ and $\gamma=C_p^{\text{mix}}/C_v^{\text{mix}}$.}
{Do \emph{not} average the $\gamma$ values: average the heat capacities, then form $\gamma$.}
{One mole of a monatomic gas ($C_v=\tfrac32 R$) is mixed with one mole of a rigid diatomic gas
($C_v=\tfrac52 R$). The ratio $\gamma$ of the mixture is:}
{$1.53$}
{$1.50$}
{$1.40$}
{$1.67$\archans{B}{$C_v^{\text{mix}}=\tfrac12\left(\tfrac32R+\tfrac52R\right)=2R$ and
$C_p^{\text{mix}}=3R$, so $\gamma=1.5$. (A) averages $\tfrac53$ and $\tfrac75$ directly; (C) and (D) are
the values for the pure diatomic and pure monatomic gases.}}

\archq{Degrees of freedom from a heat capacity}
{``A gas has $C_p$ (or $C_v$, or $\gamma$) equal to a given value; what kind of molecule is it?''}
{Invert $C_v=\tfrac f2R$, $C_p=(\tfrac f2+1)R$ or $\gamma=1+\tfrac2f$ for $f$, then match $f$ to a shape:
3 monatomic, 5 rigid linear, 6 rigid non-linear, and 2 more for each active vibration.}
{Confusing $C_p$ with $C_v$: they differ by $R$, which shifts $f$ by 2.}
{The molar heat capacity of an ideal gas at constant \emph{pressure} is $4R$. Its molecules are most likely:}
{monatomic}
{rigid and non-linear polyatomic}
{rigid diatomic}
{diatomic and vibrating\archans{B}{$(\tfrac f2+1)R=4R$ gives $f=6$: three translations and three rotations,
a rigid non-linear molecule such as water vapour. (C) treats $4R$ as close to $C_p=\tfrac72R$; (D) reads
$4R$ as $C_v$, giving $f=8$ and guessing the nearest case ($f=7$); a monatomic gas, (A), has
$C_p=\tfrac52R$.}}

\archq{Mean free path under changed conditions}
{``The mean free path is $\lambda$ under some conditions; find it when temperature or pressure changes.''}
{$\lambda=\dfrac{k_BT}{\sqrt2\,\pi d^2 P}=\dfrac{1}{\sqrt2\,\pi d^2n_V}$: $\lambda\propto T$ at fixed $P$,
$\lambda\propto1/P$ at fixed $T$, and unchanged at fixed $n_V$ (a sealed rigid vessel).}
{Decide first what is held fixed. Heating a sealed rigid box does \emph{not} change $\lambda$; heating at
constant pressure does.}
{At \SI{273}{K} a gas has a mean free path of $1500d$, where $d$ is the molecular diameter. If the pressure
is held constant and the temperature is raised to \SI{373}{K}, the new mean free path is about:}
{$1500d$}
{$1098d$}
{$1753d$}
{$2049d$\archans{D}{At constant pressure $\lambda\propto T$, so $\lambda'=1500d\times\dfrac{373}{273}
\approx2049d$. (A) treats the vessel as sealed; (B) inverts the ratio; (C) takes $\lambda\propto\sqrt T$,
the scaling of the speed rather than of the free path.}}

% ================================================================
\section{Practice Problems}\label{sec:ktpractice}

The problems below are of standard examination difficulty and are meant to be worked without looking
back at the chapter. Answers, with brief working, follow the list. Take
$R = \SI{8.314}{J.mol^{-1}.K^{-1}}$ and $k_B=\SI{1.38e-23}{J/K}$ unless stated otherwise.

\begin{enumerate}[itemsep=3pt]
\item Find $v_{\text{rms}}$ for oxygen ($M=\SI{32}{g/mol}$) at \SI{300}{K}.
\item At what temperature is the rms speed of a gas double its value at \SI{300}{K}?
\item Find the total internal energy of 2 moles of a rigid diatomic gas at \SI{300}{K}.
\item At the same temperature, by what factor is the rms speed of helium ($M=\SI{4}{g/mol}$) greater than
      that of nitrogen ($M=\SI{28}{g/mol}$)?
\item A gas has $\gamma = 4/3$. How many degrees of freedom does each molecule have, and what kind of
      molecule is it likely to be?
\item A sealed rigid container of gas is heated from \SI{300}{K} to \SI{600}{K}. What happens to
      (a) the pressure, (b) the mean free path, (c) the collision frequency?
\item One mole of a monatomic gas is mixed with one mole of a rigid diatomic gas. Find $C_v$ of the
      mixture, and its $\gamma$.
\item \emph{Conceptual.} Two different gases are at the same temperature. What is necessarily equal
      between them, and what is necessarily different (assuming different molar masses)?
\item For oxygen ($M=\SI{32}{g/mol}$) at \SI{300}{K}, find the most probable speed $v_p$, the mean
      speed $\bar v$ and the rms speed $v_{\text{rms}}$, and verify that they are in the ratio
      $\sqrt2:\sqrt{8/\pi}:\sqrt3$.
\item A gas is heated at constant volume so that its rms speed doubles. By what factor does the
      absolute temperature change, and by what factor the pressure?
\item Carbon dioxide (CO$_2$) is a \emph{linear} triatomic molecule; water vapour (H$_2$O) is
      \emph{non-linear}. Assuming rigid molecules, find $f$, $C_v$ and $\gamma$ for each, and explain
      the difference.
\item \emph{Conceptual.} Sketch how the Maxwell--Boltzmann curve changes when a gas is heated. Which way
      does the peak move, and what happens to the fast tail? Name one everyday process governed by that
      tail.
\item A vessel of volume \SI{0.050}{m^3} holds \SI{1}{mol} of nitrogen and \SI{2}{mol} of helium at
      \SI{300}{K}. Find the partial pressure of each gas and the total pressure.
\item Find the average translational kinetic energy of a gas molecule at \SI{27}{\celsius}, in joules
      and in electronvolts ($\SI{1}{eV}=\SI{1.60e-19}{J}$).
\item An ultra-high-vacuum chamber is pumped down to \SI{1e-8}{Pa} at \SI{300}{K}. Find the number of
      molecules per cubic centimetre and the mean free path, taking $d=\SI{3.7e-10}{m}$.
\end{enumerate}

\subsection*{Answers}
{\small
\begin{enumerate}[itemsep=2pt]
\item $v_{\text{rms}} = \sqrt{3(8.314)(300)/0.032} \approx \SI{484}{m/s}$.
\item $v_{\text{rms}} \propto \sqrt{T}$, so doubling the speed needs four times the absolute temperature:
      $T = \SI{1200}{K}$ (\SI{927}{\celsius}).
\item $U = n\tfrac{f}{2}RT = 2\times\tfrac52\times8.314\times300 \approx \SI{12.5}{kJ}$.
\item $v \propto 1/\sqrt{M}$: the ratio is $\sqrt{28/4} = \sqrt7 \approx 2.65$.
\item $\gamma = 1 + 2/f = 4/3$ gives $f = 6$: a rigid non-linear polyatomic molecule (3 translations
      $+$ 3 rotations), such as methane or water vapour.
\item (a) $P = n_Vk_BT$ with $n_V$ fixed: the pressure doubles. (b) $\lambda = 1/(\sqrt2\pi d^2 n_V)$
      depends only on $n_V$ and $d$, both fixed: \textbf{unchanged}. (c) $z=\bar v/\lambda \propto
      \sqrt{T}$: it rises by a factor $\sqrt2\approx1.41$.
\item $C_v^{\text{mix}} = \dfrac{1\cdot\tfrac32R + 1\cdot\tfrac52R}{2} = 2R \approx
      \SI{16.6}{J.mol^{-1}.K^{-1}}$. Then $C_p = 3R$ and $\gamma = 3/2$.
\item Necessarily equal: the average translational kinetic energy per molecule, $\tfrac32k_BT$.
      Necessarily different: the rms speed (and the other two characteristic speeds). The lighter gas
      moves faster in exactly the ratio $\sqrt{M_2/M_1}$, so that $\tfrac12 m\overline{v^2}$ still
      matches.
\item $\sqrt{RT/M} = \sqrt{8.314(300)/0.032} \approx \SI{279}{m/s}$, so
      $v_p=\sqrt2\times279 \approx \SI{395}{m/s}$, $\bar v = \sqrt{8/\pi}\times279\approx\SI{446}{m/s}$ and
      $v_{\text{rms}} = \sqrt3\times279\approx\SI{484}{m/s}$. The ratio $395:446:484 = 1:1.13:1.22$, which
      is $\sqrt2:\sqrt{8/\pi}:\sqrt3$.
\item $v_{\text{rms}}\propto\sqrt T$, so doubling it needs $T$ to rise by a factor $2^2 = 4$. At fixed
      volume $P\propto T$, so the pressure also rises by a factor 4.
\item \textbf{CO$_2$ (linear):} 3 translations $+$ 2 rotations, $f=5$, so $C_v=\tfrac52R$ and
      $\gamma = 1+\tfrac25 = 1.4$. \textbf{H$_2$O (non-linear):} 3 translations $+$ 3 rotations, $f=6$, so
      $C_v = 3R$ and $\gamma = 1+\tfrac26 = \tfrac43 \approx 1.33$. A linear molecule has no rotation about
      its own axis that can take up energy, so it has one rotational degree of freedom fewer than a bent
      one, and a larger $\gamma$. (Real CO$_2$ already has partly excited vibrations at room temperature,
      so its measured $\gamma$ is about 1.29.)
\item The peak moves \emph{right} ($v_p\propto\sqrt T$) and becomes \emph{lower and broader}, since the
      area under the curve, the total fraction, stays equal to one. The fast tail grows much fatter
      (Figure~\ref{fig:ktmaxwell}). Everyday consequences: faster evaporation, faster cooking and faster
      chemical reactions, all driven by the enlarged high-energy tail.
\item $P_i=n_iRT/V$: helium $2(8.314)(300)/0.050\approx\SI{9.98e4}{Pa}$, nitrogen
      $\approx\SI{4.99e4}{Pa}$; total (Dalton) $\approx\SI{1.50e5}{Pa}$. The heavier nitrogen exerts
      half the pressure of the helium only because there is half as much of it: pressure counts molecules,
      not mass.
\item $\overline K=\tfrac32k_BT=1.5(1.38\times10^{-23})(300)\approx\SI{6.2e-21}{J}\approx\SI{0.039}{eV}$.
\item $n_V=P/k_BT=10^{-8}/(1.38\times10^{-23}\times300)\approx\SI{2.4e12}{m^{-3}}$, that is,
      about $2.4$ million molecules per cubic centimetre, still a great many.
      $\lambda=1/(\sqrt2\pi d^2n_V)\approx\SI{6.8e5}{m}$, nearly \SI{700}{km}: a molecule crosses the
      chamber thousands of times, bouncing between the walls, before it meets another molecule.
\end{enumerate}

}

% ================================================================
\section{Challenge Problems}\label{sec:ktchallenge}

The five problems below are original problems in the spirit of olympiad and JEE Advanced collections.
The first three extend the chapter's main results; the last two apply effusion and the mean free path to
real technology. Work each before reading the answer.

\begin{challengebox}{Three speeds of a gas (extends the Maxwell distribution)}
For nitrogen ($M=\SI{28}{g/mol}$) at \SI{300}{K}, compute the rms, mean and most probable molecular speeds,
and show that their ratio is universal (independent of the gas and the temperature).
\challengeans{With $R=\SI{8.314}{J.mol^{-1}.K^{-1}}$:
\[
v_{\text{rms}}=\sqrt{\frac{3RT}{M}}\approx\SI{517}{m/s},\quad
\bar v=\sqrt{\frac{8RT}{\pi M}}\approx\SI{476}{m/s},\quad
v_p=\sqrt{\frac{2RT}{M}}\approx\SI{422}{m/s}.
\]
Their ratio is $v_{\text{rms}}:\bar v:v_p=\sqrt3:\sqrt{8/\pi}:\sqrt2\approx1.22:1.13:1$, the \emph{same
for every gas at every temperature}, because all three scale as $\sqrt{T/M}$; only the numerical constant
from the Maxwell distribution differs. Always $v_p<\bar v<v_{\text{rms}}$, since the high-speed tail pulls
the mean and the rms above the peak.}
\end{challengebox}

\begin{challengebox}{How far between collisions (extends the mean free path)}
Estimate the mean free path of an air molecule at \SI{300}{K} and \SI{1.0e5}{Pa} (molecular diameter
$d\approx\SI{3.7e-10}{m}$), and compare it with the molecular size and spacing.
\challengeans{
\[
\lambda=\frac{k_BT}{\sqrt2\,\pi d^2 P}=\frac{(1.38\times10^{-23})(300)}{\sqrt2\,\pi(3.7\times10^{-10})^2(1.0\times10^5)}
\approx\SI{68}{nm}.
\]
That is about 180 molecular diameters. The spacing between molecules is only
$n_V^{-1/3}=(k_BT/P)^{1/3}\approx\SI{3.5}{nm}$, so a molecule passes some twenty neighbours before
actually striking one. Compared with Example~\ref{ex:ktmfp}, the free path is a little longer because
the gas is warmer and slightly less dense ($\lambda\propto T/P$).}
\end{challengebox}

\begin{challengebox}{The ratio $\gamma$ of a gas mixture (extends equipartition)}
Equal numbers of moles of helium (monatomic) and nitrogen (diatomic) are mixed. Find the molar heat
capacities $C_v$ and $C_p$ and the ratio $\gamma$ of the mixture.
\challengeans{Heat capacities add per mole. For the equimolar mixture,
\[
C_v=\frac{1}{2}\left(\frac{3R}{2}+\frac{5R}{2}\right)=2R,\qquad C_p=C_v+R=3R,
\]
so $\gamma=C_p/C_v=1.5$. It lies between helium's $\tfrac53\approx1.67$ and nitrogen's $\tfrac75=1.4$, as
it must. Measuring $\gamma$ (for example from the speed of sound) is a classic way to \emph{determine} a
gas's atomicity or a mixture's composition. Note that $1.5$ is not the average of $1.67$ and $1.4$: it is
the heat capacities, not their ratios, that add.}
\end{challengebox}

\begin{challengebox}{Escaping through a pinhole (extends effusion)}
A vessel contains a mixture of hydrogen and oxygen. Gas leaks slowly through a tiny hole. By what factor
does hydrogen escape faster than oxygen, and what happens to the composition of the gas left behind? Then
apply the same idea to uranium hexafluoride, $^{235}$UF$_6$ ($M=\SI{349}{g/mol}$) and $^{238}$UF$_6$
($M=\SI{352}{g/mol}$): what is the enrichment factor of one effusion stage, and roughly how many ideal
stages raise the $^{235}$U fraction from its natural \SI{0.72}{\percent} to the \SI{4}{\percent} used in
power reactors?
\challengeans{The effusion rate is $\tfrac14 n_V\bar vA\propto 1/\sqrt M$ per molecule present (Graham's
law), so
\[
\frac{\text{rate}_{\text{H}_2}}{\text{rate}_{\text{O}_2}}=\sqrt{\frac{M_{\text{O}_2}}{M_{\text{H}_2}}}
=\sqrt{\frac{32}{2}}=4.
\]
Hydrogen escapes four times as fast, so the remaining gas becomes progressively \emph{richer in oxygen}.
For UF$_6$ the factor per stage is only $\sqrt{352/349}\approx1.0043$: each stage multiplies the ratio of
$^{235}$U to $^{238}$U by 1.0043. The ratio must grow from $0.0072/0.9928$ to $0.04/0.96$, a factor of
about 5.7, which needs $\ln5.7/\ln1.0043\approx400$ ideal stages. Real diffusion plants, with imperfect
barriers, used well over a thousand. The nuclear physics of why $^{235}$U is the useful isotope is the
subject of \chref{ch:nucleus} and the chapters after it.}
\end{challengebox}

\begin{challengebox}{When is a vacuum a vacuum? (extends the mean free path)}
In a vacuum tube or an electron-beam instrument, molecules (and electrons) must cross the apparatus,
about \SI{10}{cm}, without colliding. Find the pressure at \SI{300}{K} at which the mean free path of air
molecules ($d=\SI{3.7e-10}{m}$) equals \SI{10}{cm}, and the fraction of atmospheric pressure that this
represents.
\challengeans{From $\lambda=k_BT/\sqrt2\pi d^2P$,
\[
P=\frac{k_BT}{\sqrt2\,\pi d^2\lambda}=\frac{(1.38\times10^{-23})(300)}{\sqrt2\,\pi(3.7\times10^{-10})^2(0.10)}
\approx\SI{0.07}{Pa},
\]
less than a millionth of atmospheric pressure ($0.068/1.0\times10^{5}\approx7\times10^{-7}$). Below this
pressure the gas stops behaving like a fluid inside the apparatus: molecules travel from wall to wall,
and ideas such as viscosity and convection from \chref{ch:fluids} no longer apply. This is the regime of
effusion and of molecular beams.}
\end{challengebox}

% ================================================================
\section*{Looking Back on Part C}
\addcontentsline{toc}{section}{Looking Back on Part C}

\textbf{Part C found the laws of heat and the molecules behind them.}
\begin{itemize}
\item \textbf{Thermodynamics (\chref{ch:thermo}).} Energy is conserved once heat is counted:
$\Delta U=Q-W$, with internal energy a state function and heat and work not. Applied to the standard
processes (isothermal, adiabatic, isobaric, isochoric, and the polytropic family that contains them all),
the first law fixes what a gas does when it is squeezed, heated or allowed to expand, and why $C_p>C_v$.
The second law says that not all heat can become work: Carnot's limit $\eta_{\text{Carnot}}=1-T_c/T_h$ caps every
engine, heat flows of itself only from hot to cold, and a refrigerator must pay in work to reverse that flow. Entropy, a state variable built from the Carnot cycle, turns all of this into one statement: the entropy of an isolated system never decreases.
\item \textbf{Kinetic theory (\chref{ch:kinetic}).} All of that is molecules obeying Newton's laws.
Pressure is the momentum molecules deliver to the walls, $P=\tfrac13\rho\overline{v^2}$; temperature is
the average translational kinetic energy, $\tfrac32k_BT$ per molecule; the speeds follow the
Maxwell--Boltzmann distribution, with $v_p<\bar v<v_{\text{rms}}$; equipartition gives $\tfrac12k_BT$ to
every squared term and so predicts $C_v$, $C_p$ and $\gamma$; and the mean free path,
$1/\sqrt2\pi d^2n_V$, sets how far a molecule travels between collisions.
\end{itemize}
\textbf{What kinetic theory added to thermodynamics.} Thermodynamics took two facts on trust: the
ideal-gas law, and the claim that an ideal gas's internal energy depends on its temperature alone. Kinetic
theory derived both, and in doing so turned $C_v$ and $\gamma$ from measured numbers into predictions of
molecular shape. It also showed where the laws stop: real gases liquefy when molecules are crowded, and
equipartition fails when quantum steps are large compared with $k_BT$. The two descriptions are one
subject seen from two distances: thermodynamics from the outside, valid whatever the molecules are doing,
and kinetic theory from the inside, explaining why.

% ================================================================
\section*{Looking Back on the Volume}
\addcontentsline{toc}{section}{Looking Back on the Volume}

\textbf{Volume~II asked how bulk matter responds to force and to heat, and found that both answers come
from molecules in motion.}
\begin{itemize}
\item \textbf{Part~A: mechanical properties of matter (\chsref{ch:solids}{ch:fluids}).} A solid resists
deformation through moduli that measure the stiffness of the springs between its atoms, and stores
elastic energy exactly as a spring does in Volume~I. A fluid at rest pushes equally in all directions
with a pressure that grows with depth, which explains buoyancy and hydraulics; a moving fluid obeys
conservation of mass (continuity) and of energy (Bernoulli), and its viscosity and surface tension are the bulk signs of molecules attracting their neighbours and carrying momentum between layers.
\item \textbf{Part~B: heat and temperature (\chref{ch:thermal}).} Temperature is what thermometers agree
on in equilibrium; heat is energy in transit because of a temperature difference. Heating makes matter
expand, raises its temperature in proportion to its specific heat, or changes its phase at fixed
temperature, and heat moves by conduction, convection and radiation.
\item \textbf{Part~C: thermodynamics and the molecular picture (\chsref{ch:thermo}{ch:kinetic}).} The laws
of thermodynamics organise heat and work, and kinetic theory explains them from the motion of molecules.
\end{itemize}

\medskip
\textbf{The volume as a whole.} Notice how little was ever really separate. A modulus of elasticity, the pressure of a gas and the viscosity of a fluid are averages of molecular forces and molecular momentum; the specific heat of a solid is $3R$ per mole because each atom is a little oscillator with six
squared terms; the expansion of a heated rod and the pressure of a heated gas are both molecules moving
more vigorously; and heat itself is the energy of molecular motion in transit. Volume~I's mechanics,
applied to $10^{23}$ particles at once and read through averages, became the physics of matter: one quantity, $\tfrac12 k_BT$ per degree of freedom, seen in the pressure of a gas and in the heat capacities of gases and solids alike. The same
two ideas, energy conservation and the Boltzmann factor, will keep reappearing in later physics wherever
many particles share energy.

\medskip
\textbf{Ahead to Volume~III.} This volume treated atoms as hard, indivisible balls whose energies are a
few hundredths of an electronvolt: $k_BT\approx\SI{0.026}{eV}$ at room temperature. Volume~III goes inside
the atom to its nucleus (\chref{ch:nucleus}), where the energies are millions of electronvolts, tens of
millions of times larger, and the rules change: nuclei can decay, spontaneously and at random, and can be
split or fused. Two ideas from this volume travel with us. The first is statistics: just as the pressure
of a gas is steady although each molecule is unpredictable, the decay of a large sample of nuclei follows a
precise law although no single decay can be foreseen. The second is the fast tail of the Maxwell
distribution: in the core of the Sun, at about \SI{1.5e7}{K}, the typical thermal energy $k_BT$ is only
about \SI{1.3}{keV}, far too little to push two protons together against their electric repulsion, yet
the few nuclei in the extreme tail, helped by quantum tunnelling, fuse often enough to power the star.
